\section{Discussion}
The same trained agent expresses two qualitatively different responses to an invalid cue depending on stimulus load. At set size four, an invalid cue costs the model one frame of attentional reallocation: the spotlight, initially committed to the cued location, re-points to the uncued change one frame late and then locks on it fully (excess change-lock $+0.70$ over the uniform baseline, against $+0.75$ for a valid cue), and the behavioural psychometric confirms that uncued changes are detected, only at higher orientation magnitude (a rightward threshold shift, e.g. $0.425$ versus $0.584$ detection at $16^\circ$, with both conditions recovering to ceiling by $44^\circ$). At set size nine the reallocation does not occur at all: the model places no excess attention on the uncued change at either frame ($\approx 0$ / slightly negative), even though valid-cue locking is strong ($+0.69$ rising to $+0.86$) and the cue-frame read-out is intact. The invalid-cue cost has changed in kind, not merely in degree, becoming an effective blindness to uncued changes.

\subsection{A biased-competition / normalization mechanism}
This set-size scaling is what biased-competition and normalization accounts predict \citep{desimone1995neural,reynolds1999competitive,reynolds2000attention,reynolds2009normalization}. With more competitors in the suppressive field, the top-down bias toward the cued location is harder to overturn: reallocating attention away from a strongly favoured cued item toward an uncued event requires overcoming a deeper pool of competition, and at set size nine that reallocation fails outright. The same logic frames the effect as a load-graded extension of the classical Posner invalid-cost \citep{posner1980attention,posner1994visual}, where the cost of a misdirected cue grows with the number of competing locations \citep{carrasco2011visual}. A priority-map reading is consistent: the model commits to the cued location, and a more strongly committed map---more competitors, stronger normalization---is correspondingly harder to override, echoing collicular and priority accounts of selective commitment \citep{krauzlis2018selective,krauzlis2013superior,zenon2012attention}.

\subsection{A top-down versus bottom-up dissociation}
The collapse is specific to the cue-gated top-down read-out. An element-wise-affine variant of the same architecture, which reads change bottom-up, does not collapse at set size nine: it continues to lock on the uncued change (invalid-cue excess $\approx +0.61$ at f5 and $+0.59$ at f6, comparable to its valid-cue excess of $+0.71/+0.66$), because any sufficiently large orientation change pulls its attention by salience regardless of the cue. The price of that immunity is a weaker cueing benefit: a reader that is captured by salience cannot express the gain that a cue-gated reader derives from prioritizing the cued location. The dissociation localizes the catastrophic invalid-cue scaling to the top-down, cue-gated channel, and identifies it as the same channel that confers the cueing benefit in the first place---selection and its cost are two faces of one mechanism \citep{sprague2018rescuing,herman2020attention}.

\subsection{A falsifiable behavioural prediction}
The abolished top-down reallocation at set size nine predicts a large behavioural invalid-cue cost there: uncued changes should be largely missed rather than merely detected at higher magnitude. We state this as a prediction. More generally, the model forecasts that the behavioural invalid-cue cost in cued change detection should grow with set size---recoverable at low load, severe at high load---a directly testable claim for covert-attention experiments in primates and humans \citep{posner1980attention,luo2019attention}.

\subsection{Limitations}
Several caveats bound these conclusions. The set-size-nine attention read-out is a snapshot of a model under continued training, taken earlier than the set-size-four read-out, so the magnitudes there may shift as training proceeds. The set-size-nine behavioural psychometric has not been measured; the large invalid-cue cost at high load is therefore a prediction to be confirmed at convergence, not a measured result, and no behavioural number is reported for that condition. Change-lock is reported throughout as excess over the relevant uniform baseline ($0.25$ at set size four, $0.111$ at set size nine) to keep the quantity comparable across set sizes. Finally, the findings rest on single trained instances of each variant; replication across seeds is needed to establish that the top-down-versus-bottom-up dissociation is robust rather than idiosyncratic.
